\documentclass[11pt]{article}
\usepackage[margin=1in]{geometry}
\usepackage{amsmath,amssymb,booktabs,longtable}
\usepackage[hidelinks]{hyperref}
\newcommand{\dd}{\mathrm d}
\newcommand{\eq}{\mathrm{eq}}
\title{A missing source-assisted reaction and positive pole matching\\
The next kinetic obligations of the scalar particle bridge}
\author{Zimmerman programme: CPB-SA-2026-10-16}
\date{16 October 2026}
\begin{document}
\maketitle
\begin{abstract}
The previous pair/source abundance evolution omitted a pole-free
three-to-two reaction already present in the same scalar action.
We derive its complete tree amplitude and finite-temperature MB reaction
density. At the selected 2-MeV point and $T_d=m$, its two conjugate channels
add approximately 4.95 times the earlier pair-conversion coefficient.
Exact single, double and nested source-pole residues identify the reaction
network needed before subtracting the six-point remainder. A positive
one-pole interior matching construction has a proved smooth-profile error
bound; threshold-boundary and overlapping-pole matching remain unresolved.
The old abundance curves remain correct for their stated reduced model,
not verified predictions of the complete scalar theory. Full closure,
the galaxy-response connection and the vacuum normalization remain open.
\end{abstract}

\section{Current state and exact scope}
The previous goal turn made progress by deriving and executing a reduced
pair-conversion gas model. The current checkout is now
\texttt{bd5dbb23a2fa98f11f4bd9c13f1ef29970a06458}; its intervening commits
concern independent empirical lanes. The actual local kinetic files and
raw results were inspected before continuing. No same-route live job was
found in this checkpoint, and no earlier simulation was restarted.
All new computation handles have finished. Frozen model criteria and
other lanes are unchanged.

Use natural units and the scalar vertex conventions
\begin{equation}
 \mathcal L\supset-g(X^*\Phi^2+X\Phi^{*2})-\lambda|\Phi|^4,
 \quad q_\Phi=1,\quad q_X=2,
 \quad V_3=-2ig,\quad V_4=-4i\lambda.
\end{equation}
Here $\lambda=\lambda_{\rm tot}=\lambda_{\rm rep}+g^2/\Delta$,
$\Delta=M^2-4m^2>0$. The rate benchmark is the separately specified
$m=2$ MeV, $u_*=.02$, $r_\lambda=2$ mass-family candidate.
The equations below do not identify it with the static gravitational
polarization field. The source is treated as a narrow on-shell quasiparticle;
medium and Bose corrections are not included.

\section{The single-source cuts require more than a pair reaction}
The same action permits the two physically distinct crossings
\begin{align}
 X+\Phi^*+\Phi^*&\rightleftarrows\Phi+\Phi^*,\label{eq:SAI}\\
 X+\Phi+\Phi^*&\rightleftarrows\Phi+\Phi.\label{eq:SAJ}
\end{align}
Their conjugates are separate charge-related reactions. The first has
two identical initial daughters; the second has two identical final
daughters. Each needs one factorial $2!$ in its invariant event rate.
Neither was included in the prior pair/source evolution.

At the tree-amplitude level, cut an incoming same-charge pair from the
six-daughter reaction $2\Phi+2\Phi^*\to\Phi+\Phi^*$. With
$D_+=(p_1+p_2)^2-M^2$,
\begin{equation}
 \operatorname*{Res}_{D_+=0}\mathcal M_6=2g\mathcal M_5,
\end{equation}
where the right side is the physical source-assisted amplitude in
equation~\eqref{eq:SAI}. Cutting both incoming same-charge pairs gives
\begin{equation}
 \operatorname*{Res}_{D_+=D_-=0}\mathcal M_6
 =4g^2\mathcal M_{XX^*\to\Phi\Phi^*}.
\end{equation}
The nested pole of $\mathcal M_5$ itself has residue
$2g\mathcal M_{XX^*\to\Phi\Phi^*}$. A single real-source cut of the crossed
six-daughter reaction $3\Phi+\Phi^*\to2\Phi$ instead gives
equation~\eqref{eq:SAJ}. These identities follow by removing the cut tree
propagator and its cubic end vertex, with the stated amplitude sign convention.

The exact checks retain every labeled diagram. There are 36 six-point
trees and six one-source/four-daughter five-point trees. Single, double,
nested and crossed residues agree as exact polynomials in $g/m$ at nine
rational parameter points. External momenta are exactly on shell by using
rational spatial squared weights, rather than floating square roots.
The evidence is \texttt{source\_assisted\_matching\_run\_002}; its 72
exact controls support the implementation. The general factorization
argument is the tree cut, not a conclusion extrapolated from nine samples.
This is self-review.

The consequence for subtraction is precise: deleting only the double-source
pole and donating it to $XX^*\leftrightarrow\Phi\Phi^*$ leaves single-source
histories unmatched. The source-assisted collision channels must be supplied
before the residual six-point rate can be called fully matched.

\section{Complete five-point tree amplitude}
Take all external momenta incoming. Let $Q$ be the charge-two source,
$p$ the unique positive-charge daughter momentum, and
$r_0,r_1,r_2$ the three negative-charge daughter momenta. Then
\begin{equation}
 Q+p+r_0+r_1+r_2=0.
\end{equation}
Three trees have one quartic and one cubic vertex; three have three cubic
vertices. Summing them gives the compact exact expression
\begin{equation}
 \boxed{\mathcal M_5=-8g\sum_{k=0}^2
 \frac{\lambda+g^2/[(r_i+r_j)^2-M^2]}
 {(Q+r_k)^2-m^2},\qquad \{i,j,k\}=\{0,1,2\}.}
\label{eq:SAamp}
\end{equation}
For crossing~\eqref{eq:SAI}, one source denominator can be on shell:
the invariant mass of the two incoming negative daughters can equal $M$.
Its residue is the previously stated pair-scattering amplitude. A full
thermal rate for this crossing is therefore not obtained by simply squaring
the zero-width expression.

\section{The other crossing has no internal real-particle pole}
For equation~\eqref{eq:SAJ}, let $Q,p,r$ be future-directed incoming
source, positive daughter and negative daughter; let $k,l$ be the two
future-directed outgoing positive daughters. In equation~\eqref{eq:SAamp},
the negative momenta are $r,-k,-l$.

The source denominators obey
\begin{align}
 (k+l)^2-M^2&=s-M^2\ge4m(M+m)>0,\\
 (r-k)^2-M^2&\le-M^2,\qquad
 (r-l)^2-M^2\le-M^2.
\end{align}
For the daughter denominators,
\begin{equation}
 (Q+r)^2-m^2\ge M(M+2m)>0.
\end{equation}
The vector $K=Q-k=l-(p+r)$ is strictly spacelike. If $K$ were future
causal, $l=(p+r)+K$ would have mass at least $2m$, impossible. If $K$
were past causal, $k=Q-K$ would have mass at least $M>m$, also impossible.
$K=0$ would identify unequal on-shell masses. Thus
\begin{equation}
 (Q-k)^2-m^2<-m^2,\qquad (Q-l)^2-m^2<-m^2.
\end{equation}
The whole tree is pole-free throughout the physical region. This is a
uniform kinematic statement, not just a numerical absence of poles.

Since $\lambda>g^2/M^2$, an explicit uniform amplitude bound is
\begin{equation}
 |\mathcal M_5|\le C_5\equiv8|g|
 \left[\frac{\lambda+g^2/[4m(M+m)]}{M(M+2m)}
 +\frac{2\lambda}{m^2}\right].
\end{equation}
It also gives a finite positive analytic bound on the MB reaction density:
\begin{equation}
 \gamma_3^\eq\le\frac{C_5^2}{16\pi}I_X I_\Phi^2,
 \qquad I_i=\int\frac{\dd^3p}{(2\pi)^3 2E_i}e^{-E_i/T_d}
 =\frac{m_iT_d}{4\pi^2}K_1(m_i/T_d).
\end{equation}
The factor $1/(16\pi)$ is the maximal two-body phase space with its
identical-final $1/2!$. The computed rates respect this bound. It is loose
and does not replace quadrature-error assessment.

\section{Threshold amplitude and event normalization}
For all three incoming particles at rest in crossing~\eqref{eq:SAJ},
the final energies are $(M+2m)/2$. Evaluation of equation~\eqref{eq:SAamp}
gives
\begin{equation}
 \mathcal M_{5,0}=
 \frac{8g\lambda(M+m)}{Mm(M+2m)}
 -\frac{2g^3(5M+8m)}{mM^2(M+m)(M+2m)}.
\end{equation}
It is nonzero for the target range $2m<M<4m$, $g>0$,
$\lambda_{\rm rep}>0$. To see the sign, factor out
$8g/[Mm(M+2m)]$ and substitute
$\lambda=\lambda_{\rm rep}+g^2/(M^2-4m^2)$.
The remaining $g^2$ numerator is proportional to
\begin{equation}
 -x^3+24x+32=64+(4-x)(x^2+4x-8)>0,\qquad x=M/m\in(2,4).
\end{equation}
The $\lambda_{\rm rep}$ part is positive too. This establishes the
nonvanishing vacuum-tree threshold amplitude on that stated domain.

For a dilute nonrelativistic uncorrelated gas, the forward event rate of
one crossing~\eqref{eq:SAJ} is
\begin{equation}
 R_3=K_{3,0}n_Xn_\Phi n_{\Phi^*},\qquad
 K_{3,0}=\frac{\sqrt{1-4m^2/(M+2m)^2}\,|\mathcal M_{5,0}|^2}
 {128\pi M m^2}.
\end{equation}
The initial invariant measures give $1/[(2M)(2m)^2]$, and the final phase
space with the identical factor gives $\beta_f/(16\pi)$. Finite-temperature
amplitudes cannot be replaced by this threshold coefficient in a warm gas.

\section{A resolved finite-temperature kernel}
With classical equilibrium distributions and the final $1/2!$ included,
\begin{align}
 \gamma_3^\eq(T_d)&=\frac{T_d}{16\pi^3}
 \int_{(M+2m)^2}^\infty\sqrt{s}K_1(\sqrt{s}/T_d)\mathcal I_5(s)\dd s,\\
 \mathcal I_5(s)&=\frac{\beta_f}{1024\pi^4s}
 \int_{4m^2}^{(\sqrt{s}-M)^2}\!
 \sqrt{\mathcal K(s,M^2,u)}\sqrt{1-4m^2/u}
 \langle|\mathcal M_5|^2\rangle_{\rm angles}\dd u.
\end{align}
Here $\mathcal K$ is the K\"all\'en polynomial, $\beta_f=\sqrt{1-4m^2/s}$,
and the angular brackets are a normalized average over the initial relative
polar angle and final polar/relative azimuthal angles. The three-body phase
space follows by inserting the pair invariant $u$ and two two-body phase
spaces. The thermal total-momentum integral gives
$T_d\sqrt{s}K_1(\sqrt{s}/T_d)/(8\pi^3)$; the additional factor $1/2!$
produces the displayed coefficient. The invariant rate conventions agree
with the primary-source discussion in \cite{SAthermal}.

The code uses generalized Laguerre quadrature for total kinetic energy,
Jacobi quadrature for the initial Dalitz invariant, and polar Gaussian and
periodic azimuthal quadratures. Exponentially scaled Bessels avoid
underflow in $\widehat\gamma_3=\gamma_3^\eq e^{4m/T_d}$.
All internal denominators retain their proved physical signs.

The original coarse run failed its 2-percent refinement gate at $T_d=m$.
A separate resolution experiment found that refining energy/Dalitz nodes
alone changed the warm values negligibly, while the angular grid mattered.
The failed run remains intact. The authoritative refined run increases
angular resolution and retains the original gate. At $T_d=3m$, its two
last angular grids differ by .261 percent. These differences are finite
numerical evidence, not rigorous error bars. The low-$T_d/m=.003$ rate
approaches the threshold-density normalization within 2 percent, with
remaining finite-temperature corrections retained.

For source partial equilibrium $z_X=z_\Phi^2$ and equal charge populations,
the two conjugate source-assisted reactions add
\begin{equation}
 P_3=2\gamma_3^\eq(z_\Phi^2-z_\Phi^4)
\end{equation}
to $\dot w+3Hw$. The earlier pair contribution is
$P_2=\gamma_2^\eq(z_\Phi^2-z_\Phi^4)$. Consequently their coefficient
ratio is independent of the shared fugacity in this MB reduction:
\begin{center}
\begin{tabular}{r r r}
\toprule $T_d/m$ & $2\gamma_3^\eq/\gamma_2^\eq$ & $\gamma_3/\gamma_{3,\rm NR}$\\\midrule
.003 & .0017634 & .98001\\
.01 & .0096558 & .93529\\
.03 & .047251 & .82095\\
.1 & .25802 & .53614\\
.3 & 1.10193 & .19128\\
1 & 4.95265 & .016137\\
3 & 20.7964 & .00039732\\\bottomrule
\end{tabular}
\end{center}
In particular this omitted reaction is not justified as negligible at the
warmer temperatures used in the prior reduced abundance comparison. Its
threshold approximation would greatly overestimate the warm reaction.
These are vacuum-tree, common-MB coefficients; they are not observed rates
or a new primordial abundance calculation.

The exact residue and new kernel computations also explain why the prior
yield ratio 1.97924 must remain a result for its pair-only network. No
updated abundance ratio is asserted until the source-assisted channels,
remaining subtractions and applicable quantum/medium corrections have
been treated in the same evolution.

\section{A positive interior pole matching lemma}
The general real-intermediate-state problem is known \cite{SARIS}. The
following is a scoped construction for a single pole strictly inside the
physical integration domain. It is not a completed thermal prescription.

Write $D=s-M^2$, $a=M\Gamma>0$ and
\begin{equation}
 \mathcal A(D)=B(D)+\frac{R(D)}{D+ia},\qquad -W\le D\le W.
\end{equation}
Assume real $B,R$, nonnegative integration weight $v$, and define
$F=vR^2\in C^2$, $G=vRB\in C^1$ on the symmetric window. Let
$|F''|\le F_2$, $|G'|\le G_1$. Keep the full nonnegative squared
amplitude outside the window. Inside it keep the background $vB^2$ and
replace the resonant contribution by the positive local pole term
\begin{equation}
 \frac{2F(0)}a\arctan(W/a)
 =\frac{\pi F(0)}a Z_W,\qquad
 Z_W=\frac2\pi\arctan(W/a)\in(0,1).
\end{equation}
Equivalently the scattering amplitude uses the pole propagator only
outside the window, with the retained real-state spectral weight counted
separately. All retained squared-amplitude and pole terms are nonnegative.

The difference from the original full window integral is
\begin{equation}
 E_W=\int_{-W}^W\frac{F(D)-F(0)}{D^2+a^2}\dd D
 +2\int_{-W}^W\frac{G(D)D}{D^2+a^2}\dd D.
\end{equation}
Odd constant/linear pieces cancel, so Taylor's theorem and the mean-value
bound prove
\begin{equation}
 \boxed{|E_W|\le(F_2+4G_1)[W-a\arctan(W/a)].}
\end{equation}
This bounds the omitted regular interference as well as variation of the
resonant coefficient. If one instead uses the full unweighted pole term
$\pi F(0)/a$, its additional known difference
$[\pi-2\arctan(W/a)]F(0)/a$ must be counted. $Z_W$ is a matching
coefficient of the window, not an independently measured matter fraction.

The numerical benchmark uses $v(D)=e^{-D/L}$ and constant real $R,B$,
with $W=\sqrt{aL}$. Its exact bound has
$F_2=R^2 e^{W/L}/L^2$, $G_1=|RB|e^{W/L}/L$.
Twenty width/background cases satisfy the bound; a direct full-window
integration supplies an independent control at a moderate width.
At the selected vacuum source point,
\begin{equation}
 \Gamma=1.05750\,10^{-11}\ \mathrm{GeV},\quad
 a=4.23084\,10^{-14}\ \mathrm{GeV}^2.
\end{equation}
Assigning $L=\Delta/4$ solely as a scale diagnostic gives
$W=8.22760\,10^{-12}$ GeV$^2$,
$W/a=194.47$, $W/\Delta=.001286$, and $Z_W=.996726$.
The width/window/threshold-gap hierarchy is available in vacuum; the
actual thermal $F_2,G_1$ have not been bounded here.

The lemma fails to cover a pole window truncated by a phase-space endpoint,
complex backgrounds with absorptive interference, and overlapping poles
without further matching. In particular the upper initial-pair invariant
can approach $M^2$ near the source-pair threshold. A fixed symmetric window
cannot simply be used there. Those physical boundary and double-pole
patches remain open, rather than being hidden in the finite benchmark.

\section{What is established and what remains}
The current evidence consists of exact source-pole factorization,
a complete pole-free source-assisted amplitude, a bounded/resolved thermal
kernel, and the interior matching lemma. It supplies a concrete way to
advance the previously unspecified collision remainder. It does not supply
the whole remainder or make the earlier abundance curves complete.

Next add this pole-free kernel to the energy/entropy-conserving evolution,
with a sufficiently resolved temperature table and its error carried
through the result. In parallel as mathematical dependencies, match
crossing~\eqref{eq:SAI} and the six-point remainder using actual thermal
profile bounds, threshold endpoints and the exact pole residues. Do not
add an unsubtracted raw six-point rate to an explicit real-source network.
The additional reaction may change the dilution/source-equilibrium gates;
their old outcomes cannot be inherited without re-evolving.

This is a collision calculation for chosen massive fields. It neither
derives the Higgs-era initial state nor identifies the fields with the
galactic polarization medium. The vacuum offset, $\kappa$, the acceleration
scale postulate, physical halo settling, finite mass supply, source
selection, relativistic lensing and native cosmology remain unresolved.
Internal rest/kinetic conversion still is not automatically transfer to
dark energy. The pressure-aware exchange identity from the previous
checkpoint remains the applicable bookkeeping condition.

Novelty boundary: source-assisted number change, tree factorization,
thermal phase-space integration and real-state matching have established
prior art. This report gives their explicit application to the current
carrier and a scoped mathematical error estimate. No global novelty,
observational identification or independent-referee claim is made.

\begin{thebibliography}{9}
\bibitem{SAthermal} V. De Romeri et al.,
\emph{Neutrino dark matter and the Higgs portal: improved freeze-in analysis},
arXiv:2003.12606v2, invariant reaction measure and flux, equations (21)--(25).
\url{https://arxiv.org/html/2003.12606v2}.
The common-MB specialization and five-point phase space here are derived
for this scalar model; its sterile-neutrino results are not imported.
\bibitem{SARIS} K. Ala-Mattinen et al.,
\emph{Anatomy of real intermediate state-subtraction scheme},
arXiv:2309.16615v2, sections I--V.
\url{https://arxiv.org/html/2309.16615v2}.
The primary source was inspected for the on-shell overlap and cut matching
framework. It does not provide our physical boundary profiles or complete
multi-source subtraction.
\end{thebibliography}
\end{document}
